% Aula 06 --- quem é sensível à escala das covariáveis, e por quê.
Um colega recebeu um banco em que a coluna \texttt{renda} está em reais e a coluna
\texttt{altura} em metros, e pergunta se precisa padronizar antes de ajustar. A
resposta honesta é "depende do método" --- e o critério que decide cabe numa
pergunta: \textbf{o método soma coisas vindas de colunas diferentes?}

Se soma, as unidades precisam ser comparáveis, e trocar uma delas estraga a
comparação. Se não soma, a unidade é irrelevante.
\begin{itens}
  \item \pts{2,4 --- 0,4 cada} Para cada método abaixo, diga se multiplicar uma
        única coluna por $1000$ \textbf{muda} ou \textbf{não muda} a predição,
        justificando em uma linha onde está --- ou onde não está --- a soma sobre
        colunas.
        \begin{enumerate}[label=(\roman*),leftmargin=1.0cm,nosep]
          \item Mínimos quadrados ordinários.
          \item Regressão Ridge com $\lambda>0$ fixo.
          \item Árvore de regressão.
          \item Floresta aleatória.
          \item KNN com $k$ fixo.
          \item \emph{Spline} de base truncada, ajustado por mínimos quadrados.
        \end{enumerate}
  \item \pts{0,7} Entre os que mudam, qual você espera que seja o mais afetado, e
        por quê? Compare o \emph{modo} como a escala entra em cada um.
  \item \pts{0,9} Um colega observa que, ao reescalar uma coluna, o risco do Lasso
        \emph{melhorou}, e conclui que não padronizar foi uma boa ideia. Explique
        o erro do raciocínio.
\end{itens}

\begin{solucao}
\textbf{(a)}

\begin{center}\small
\renewcommand{\arraystretch}{1.25}
\begin{tabular}{@{}clp{8.2cm}@{}}
\toprule
 & muda? & onde está (ou não) a soma sobre colunas \\
\midrule
(i) & \textbf{não} & O MQO é equivariante: o coeficiente da coluna se divide por
$1000$ e o produto $\beta_jx_j$ fica igual. Nada na função-objetivo compara
coeficientes entre si. \\
(ii) & \textbf{sim} & A penalidade $\lambda\sum_j\beta_j^2$ \emph{soma}
coeficientes de colunas diferentes; mudar a unidade de uma delas muda quanto ela
pesa nessa soma. \\
(iii) & \textbf{não} & A árvore só pergunta "$x_j\le t$?". O corte $t$ é
multiplicado junto, e a pergunta separa exatamente as mesmas observações. \\
(iv) & \textbf{não} & É uma média de árvores, e cada uma é indiferente; a
indiferença atravessa a agregação. \\
(v) & \textbf{sim} & A distância euclidiana $\sum_j(x_{ij}-x_{\ell j})^2$ é uma
soma sobre colunas. \\
(vi) & \textbf{não} & É um MQO sobre uma base fixa de funções de \emph{uma}
coluna; vale o mesmo argumento de (i). \\
\bottomrule
\end{tabular}
\end{center}

\smallskip
\textbf{(b)} O \textbf{KNN}, e a diferença é de ordem, não de grau.

Nos métodos penalizados a escala entra \emph{uma vez}, de forma linear, na
comparação entre coeficientes. Na distância euclidiana ela entra \textbf{ao
quadrado}: multiplicar a coluna por $1000$ multiplica as diferenças naquela
coordenada por $1000$, e a contribuição delas para a distância por $10^6$.

O efeito é que a coluna reescalada passa a decidir sozinha quem é vizinho de quem.
As demais não são apenas menos importantes --- elas somem da conta, e o método
deixa de usar a informação que carregavam.

\smallskip
\textbf{(c)} O erro é confundir um resultado com um argumento.

A melhora aconteceu porque a coluna reescalada era, naquele problema, a de maior
efeito verdadeiro. Deixá-la $1000$ vezes maior tornou o coeficiente dela $1000$
vezes menor, e um coeficiente minúsculo quase não pesa em $\sum_j\abs{\beta_j}$ ---
a coluna ficou praticamente isenta da penalidade. Isentar de penalidade justamente
a coluna que mais importa calhou de ajudar.

Calhou: é sorte, não método. Aplique a mesma reescala a uma coluna
\emph{irrelevante} e o efeito se inverte --- ela ganharia liberdade para receber um
coeficiente grande sem pagar por isso, e o ajuste pioraria.

O ponto de fundo é que, nos dois casos, quem decidiu quanto regularizar cada
coluna não foi o analista: foi a \textbf{unidade de medida} em que o dado chegou.
E é por isso que o argumento do colega não se salva nem quando o resultado o
favorece --- um procedimento cujo desempenho depende de a coluna ter sido anotada
em metros ou em milímetros está fora de controle, e ter dado certo desta vez não
muda isso.
\end{solucao}
